\subsection{Wang et al.\ (2018): DeePMD-kit, water demonstration}
\label{sec:appendix-example-wang2018dpkit}

\paragraph{Q1 -- Bibliographic identification.}
Wang, Zhang, Han, and E document DeePMD-kit, a TensorFlow-based deep
learning package for many-body potential energy representation and
molecular dynamics. Published in \emph{Comput.\ Phys.\ Commun.} 228,
178--184 (2018); arXiv:1712.03641. We exercise the protocol on the
single liquid-water demonstration in \S IV; the package's broader
multi-system, multi-method capabilities are flagged as the central
gap in Q12.
\lstinputlisting[language=turtle]{artifacts/kg/listings/wang2018dpkit/q01-bibliographic.ttl}

\paragraph{Q2 -- Material system.}
The demonstration system is liquid water (64 H$_2$O molecules under
periodic boundary conditions, ${\sim}330$\,K).
\lstinputlisting[language=turtle]{artifacts/kg/listings/wang2018dpkit/q02-material.ttl}

\paragraph{Q3 -- Reference calculation method.}
Reference data come from DFT calculations.
\lstinputlisting[language=turtle]{artifacts/kg/listings/wang2018dpkit/q03-reference-method-type.ttl}

\paragraph{Q4 -- Reference settings.}
The functional is PBE0+TS (hybrid PBE0 with Tkatchenko--Scheffler
dispersion). Other DFT settings (engine, plane-wave or basis-set
choice, k-mesh) are not stated in the demonstration section.
\lstinputlisting[language=turtle]{artifacts/kg/listings/wang2018dpkit/q04-reference-settings.ttl}

\paragraph{Q5 -- Training dataset.}
The reference set is 40{,}000 frames from a 20\,ps NVT AIMD trajectory
of 64 water molecules at 330\,K (38{,}000 train + 2{,}000 test).
Energies and forces are covered. Provenance is \emph{published}
(reused from the underlying AIMD work).
\lstinputlisting[language=turtle]{artifacts/kg/listings/wang2018dpkit/q05-dataset.ttl}

\paragraph{Q6 -- Sampling strategies.}
A single sampling strategy applies: vibrational sampling via NVT
AIMD with a 0.0005\,ps stride.
\lstinputlisting[language=turtle]{artifacts/kg/listings/wang2018dpkit/q06-sampling.ttl}

\paragraph{Q7 -- MLIP method, components, implementation.}
The method is Deep Potential Molecular Dynamics (DeePMD): a per-atom
DNN over local-frame radial+angular descriptors (1/R, x/R, y/R,
z/R) sorted by chemical species, with five hidden layers (240, 120,
60, 30, 10) and tanh activations. Loss is a weighted MSE on energy
and force (no virial) with annealed prefactors. Optimisation is Adam
with exponential learning-rate decay. The implementation is
DeePMD-kit v0.1 (Python/C++ on TensorFlow), with LAMMPS pair-style
\texttt{deepmd} and i-PI \texttt{dp\_ipi} clients shipping in the
package.
\lstinputlisting[language=turtle]{artifacts/kg/listings/wang2018dpkit/q07-method.ttl}

\paragraph{Q8 -- Hyperparameter settings.}
Reported settings: cutoff radius $r_c = 6.0$\,\AA; per-atom hidden
architecture $(240, 120, 60, 30, 10)$; initial learning rate
$10^{-3}$; learning-rate decay rate 0.95; decay applied every 5{,}000
steps.
\lstinputlisting[language=turtle]{artifacts/kg/listings/wang2018dpkit/q08-hyperparameter-settings.ttl}

\paragraph{Q9 -- MLIP run and trained model.}
A single run of $10^6$ batches at batch size 4 produces the
demonstration DeePMD water model.
\lstinputlisting[language=turtle]{artifacts/kg/listings/wang2018dpkit/q09-run-and-model.ttl}

\paragraph{Q10 -- Benchmark results.}
Test-set RMSEs: 0.028\,eV on the 64-water-cell total energy and
0.024\,eV/\AA{} on atomic-force components.
\lstinputlisting[language=turtle]{artifacts/kg/listings/wang2018dpkit/q10-benchmark-results.ttl}

\paragraph{Q11 -- Computational resources.}
The paper records training duration of approximately 16 hours on an
Intel Core i7-3770 CPU with 32\,GB RAM and 4 OpenMP threads. GPU
hours, peak memory, and per-atom inference time are not recorded.
\lstinputlisting[language=turtle]{artifacts/kg/listings/wang2018dpkit/q11-resources.ttl}

\paragraph{Q12 -- Gaps.}
The central gap is one of \emph{scope}: this is a software paper,
and the protocol's "one paper $\equiv$ one method $\times$ one material
system" framing captures only the demonstration example. The
package's multi-system support, alternate descriptor choices,
ensemble training, and downstream MD/i-PI tooling are not encoded.
The demonstration's DFT settings (engine, k-mesh, energy cutoff) are
also not surfaced beyond the functional choice; and inference-time
data are not reported even though the paper benchmarks DeePMD-LAMMPS
performance qualitatively.
